% TOP_PROBLEM: {"id": 10000061, "title": "Finite graphs whose every ball is an expander", "category_id": 3, "category": {"id": 3, "name": "graph_theory", "display_name": "Graph Theory", "description": "Problems involving graphs, networks, and their properties.", "slug": "graph-theory", "order_index": 3, "created_at": "2026-07-31T15:26:25.671Z"}, "status": "open", "problem_number": "AMR-099-0061", "set_id": 13}
% FIRST_POSTED: 2026-10-01
% ENTRY_KIND: statement_only
% UnsolvedMath ID 10000061; source catalogue code AMR-099-0061
% !TeX program = lualatex
% Definition and sources only; this file is not an LLM solution attempt.
\documentclass[11pt,a4paper]{article}
\usepackage[margin=25mm]{geometry}
\usepackage{fontspec}
\setmainfont{lmroman10-regular.otf}[BoldFont=lmroman10-bold.otf,
 ItalicFont=lmroman10-italic.otf,BoldItalicFont=lmroman10-bolditalic.otf]
\usepackage{amsmath,amssymb,mathtools}
\usepackage{unicode-math}
\setmathfont{latinmodern-math.otf}
\AtBeginDocument{\let\setminus\smallsetminus}
\usepackage{xurl,hyperref}
\hypersetup{colorlinks=true,urlcolor=blue}
\setcounter{secnumdepth}{0}
\setlength{\parindent}{0pt}
\setlength{\parskip}{5pt}
\setlength{\emergencystretch}{3em}
\begin{document}
\section*{Finite graphs whose every ball is an expander}\label{10000061}
{\small UnsolvedMath ID 10000061 · AMR-099-0061 · Graph Theory · AMR Open Problem Lists}
\subsection{Definitions and mathematical statement}
For a finite connected graph $H$ with at least two vertices, define its edge-expansion constant by
\[
h_E(H)=\min_{0<|A|\le|V(H)|/2}
\frac{|E_H(A,V(H)\setminus A)|}{|A|},
\]
where each edge crossing the indicated cut is counted once. Set $h_E(H)=+\infty$ for a singleton, so that trivial balls impose no restriction.

Does there exist an integer $d\ge2$, a constant $h>0$, and a sequence of finite connected simple $d$-regular graphs $G_n$ with $|V(G_n)|\to\infty$, such that
\[
h_E\bigl(G_n[B_{G_n}(v,r)]\bigr)\ge h
\]
for every $n$, every vertex $v\in V(G_n)$, and every integer radius $r\ge0$? Here balls are defined using the graph metric of $G_n$, and $G_n[B]$ retains exactly those original edges with both endpoints in $B$.

The source conjectures that no such family exists. The degree $d$ and expansion constant $h$ must remain fixed throughout the family and across all ball radii and centers. The induced balls are generally not regular, even though the ambient graphs are.

Taking a radius at least the diameter shows that the ambient graphs would themselves form an expander family, but that alone is far weaker than the requirement. Boundary edges leaving a ball cannot help its internal expansion. In particular, tree-like neighborhoods of increasing radius provide a natural obstruction: their large branches can be separated by very small internal cuts. The question asks whether a bounded-degree family can avoid such obstructions at every location and every scale.
\subsection{Short English statement}
Can arbitrarily large regular graphs have every induced metric ball expand by one uniform positive edge-expansion constant?
\subsection{Sources}
\begin{itemize}
\item[S1] ULAM.AI and the UnsolvedMath Contributors, supplied problems.json, record 10000061 (AMR-099-0061), AMR Open Problem Lists. Catalogue identity and source transcription; CC BY 4.0.
\url{https://huggingface.co/datasets/ulamai/UnsolvedMath}.
\item[S2] Benjamini - Euclidean vs. Graph Metric, Question 5.4, PDF page 6. Original source indexed by AMR-099-0061.
\url{https://www.wisdom.weizmann.ac.il/~itai/erd100.pdf#page=6}.
\item[S3] I. Benjamini, Euclidean vs. Graph Metric, primary numbered question and conventions.
\url{https://arquivo.pt/noFrame/replay/20201231041538id_/http://www.wisdom.weizmann.ac.il/~itai/erd100.pdf}.
\end{itemize}
\end{document}
